\begin{requirement}{HLR-SCOPE.V1-01}[draft][stakeholderRequirement]{First Functional Version}

  \localreqdiagram*{hlr05}{\reqid}{\reqdesc}

  WHEN two or more alternatives each satisfy the requirements an architecture
  decision must meet, the decision shall adopt the alternative that requires
  this project to build the fewest mechanisms of its own --- counted as
  components, public interfaces, background tasks, and on-disk structures
  built rather than adopted --- even at the cost of peak space or I/O
  efficiency.

  This specification targets a \emph{first functional version}: proven
  components doing the heavy lifting are preferred to bespoke ones. The
  requirement is verified by inspection --- each architecture decision
  record's \emph{Alternatives considered} section states what an alternative
  rejected on these grounds would have required the project to build. A later
  version of this specification may replace this requirement.

  \reqdiaglink{global}{\textit{View complete diagram.}}

  \reqcover{satisfy}{ADR}
\end{requirement}

\begin{requirement}{HLR-PERSISTENCE.DURABLE-01}[draft][stakeholderRequirement]{Durable Persistence}

  \localreqdiagram*{hlr01}{\reqid}{\reqdesc}

  WHEN the cache process exits --- in an orderly way or by crashing --- the
  cache shall serve, from the process that starts next, every cached HTTP
  object whose storage had completed before the exit.

  A process exit leaves the host's own state intact; what the cache owes
  after the host itself crashes or loses power is stated separately in
  \reqref{HLR-PERSISTENCE.CRASHSAFE-01}, and the size of the cache it can
  hold in \reqref{HLR-PERSISTENCE.CAPACITY-01}.

  \reqdiaglink{global}{\textit{View complete diagram.}}

  \reqexpect{satisfy}{ADR-PERSISTENCE.MECHANISM-01}{Establish the storage
  substrate and mechanism that persists cached objects across an exit of the
  cache process.}
  \reqcover{verify}{TC}
\end{requirement}

\begin{requirement}{HLR-PERSISTENCE.CAPACITY-01}[draft][stakeholderRequirement]{Capacity beyond RAM}

  \localreqdiagram*{hlr06}{\reqid}{\reqdesc}

  The cache shall hold a set of cached HTTP objects whose total size exceeds
  the RAM installed on its host.

  \reqdiaglink{global}{\textit{View complete diagram.}}

  \reqexpect{satisfy}{ADR-PERSISTENCE.MECHANISM-01}{Choose a storage
  substrate whose capacity is bounded by persistent storage rather than by
  process RAM.}
  \reqcover{verify}{TC}
\end{requirement}

\begin{requirement}{HLR-PERSISTENCE.CRASHSAFE-01}[draft][stakeholderRequirement]{Host-crash safety}

  \localreqdiagram*{hlr07}{\reqid}{\reqdesc}

  IF the cache's host crashes or loses power, THEN the cache, once running
  again, shall serve no object whose content differs from what was stored
  and no object whose purge it had acknowledged before the crash.

  Cached objects can be fetched again from their origin, so losing some of
  them in a host crash --- each then simply a cache miss --- is acceptable;
  serving wrong or purged content is not.

  This assumes a storage stack that honours flush requests --- what it
  acknowledges as flushed is on stable storage --- and that returns the data
  written to it or an error. Detecting data corrupted at rest, or lost by a
  device that acknowledges flushes it has not made, is the storage layer's
  responsibility, not the cache's.

  \reqdiaglink{global}{\textit{View complete diagram.}}

  \reqexpect{satisfy}{ADR-PERSISTENCE.DURABILITY-01}{Decide which writes are
  synchronised to stable storage, and what integrity check detects the rest,
  so that a host crash can lose objects but never make the cache serve
  corrupted or purged content.}
  \reqcover{verify}{TC}
\end{requirement}

\begin{requirement}{HLR-EFFICIENCY.SPACE-01}[draft][stakeholderRequirement]{Space Efficiency}

  \localreqdiagram*{hlr02}{\reqid}{\reqdesc}

  The cache shall hold the objects it caches in operation with an on-disk
  space overhead of no more than 25\,\% of those objects' own metadata and
  body bytes.

  The overhead is the storage the cache's directory tree occupies on disk ---
  allocated blocks, not file lengths --- minus the metadata and body bytes of
  the objects it holds, so it counts record headers, reserved headroom,
  block rounding, per-file and per-directory costs, and leftover files alike.
  It is measured over what the deployment actually caches, because overhead
  depends strongly on object sizes: a 500-byte object in a 4~KiB block carries
  roughly 700\,\% overhead, a 1~MiB object well under 1\,\%. The first
  functional version knowingly does not meet this requirement for workloads
  dominated by small objects: \reqref{ADR-EFFICIENCY.COMPACTION-01} defers the
  machinery that would, and declares the conflict.

  \reqdiaglink{global}{\textit{View complete diagram.}}

  \reqcover{deriveReqt}{REQ}
  \reqcover{verify}{TC}
\end{requirement}

\begin{requirement}{HLR-PERFORMANCE.LOOKUP-01}[draft][stakeholderRequirement]{Lookup latency flat with size}

  \localreqdiagram*{hlr08}{\reqid}{\reqdesc}

  WHILE the cache holds up to $10^7$ cached HTTP objects, the cache shall
  keep its 99th-percentile lookup latency no more than 25\,\% above the
  99th-percentile lookup latency it has while holding $10^4$ objects, on the
  same host under the same conditions.

  \emph{Lookup latency} is the time from the start of a cache lookup until
  the cache knows whether the request is a hit --- with the stored response
  head in hand --- or a miss; delivering the body is not part of it. ``The
  same conditions'' means the same storage device and file system, the same
  page-cache state, and the same mix of hits and misses over uniformly chosen
  keys. The bound is relative rather than absolute because absolute lookup
  latency depends mostly on the host's storage device and memory, which the
  deployment chooses; what the cache itself owes is that growing towards its
  capacity does not make every request slower. $10^7$ objects is a size a
  test host can hold; derived design targets may reach further.

  \reqdiaglink{global}{\textit{View complete diagram.}}

  \reqcover{deriveReqt}{REQ}
  \reqcover{verify}{TC}
\end{requirement}

\begin{requirement}{HLR-COMPAT.PINGORA-01}[draft][stakeholderRequirement]{Compatible with Pingora}

  \localreqdiagram*{hlr03}{\reqid}{\reqdesc}

  The cache shall integrate with the \texttt{pingora-cache} 0.9 release line
  --- version 0.9.0 and its patch releases --- through only the extension
  points those releases expose publicly, without modification of the Pingora
  source.

  The \texttt{Storage} backend is the primary integration point, but
  preserving cache semantics across restarts may require additional,
  separately-configured components --- for example eviction-state coordination
  (\reqref{FRE-COMPAT.EVICTION-01}) --- which meet this requirement as long as
  they, too, use only public extension points. Verification is by inspection
  of the deliverable's dependency manifest (Pingora crates constrained to
  0.9.x, none patched or vendored) and of the interfaces it uses. A patch
  release of the 0.9 line is taken as it comes; adopting a later minor
  release is a revision of this requirement, not an implicit obligation.
  What the persisted cache owes across any Pingora change is stated in
  \reqref{FRE-COMPAT.FORMATUPGRADE-01}. This version is such a revision: it
  moves the target from 0.8.1 to the 0.9 line, whose \texttt{Storage} and
  \texttt{EvictionManager} surfaces differ in ways the requirements below name
  explicitly.

  \reqdiaglink{global}{\textit{View complete diagram.}}

  \reqcover{deriveReqt}{FRE}
\end{requirement}

\begin{requirement}{HLR-COMPLIANCE.RFC-01}[draft][stakeholderRequirement]{No added RFC 9111/9110 deviation}

  \localreqdiagrams*{hlr04}{\reqid}{4}{\reqdesc}

  The cache's observable HTTP caching behaviour shall match that of the same
  proxy configured with the stock \texttt{MemCache} backend of
  \texttt{pingora-cache}, in every respect governed by RFC~9111 (\emph{HTTP Caching}) and
  RFC~9110 (\emph{HTTP Semantics}).

  \texttt{pingora-cache} applies these semantics --- freshness calculation,
  conditional revalidation, \texttt{Cache-Control} handling --- without
  claiming full conformance. This requirement therefore asks the backend to
  add no deviation of its own, not to remove the framework's. Behaviour
  observed across a process restart, which \texttt{MemCache} cannot exhibit,
  is compared against the behaviour \texttt{MemCache} shows within a single
  process run.

  \reqdiaglink{global}{\textit{View complete diagram.}}

  \reqcover{deriveReqt}{FRE}
  \reqcover{verify}{TC}
\end{requirement}
